% !PRINTED 148
\weberpage{148}
Der Grad $g_2$ ergiebt sich genau wie der von $Q_0$ und $Q_{0,1}$. Wenn aber
durch $Q_{0,1,2}$ noch eine vierte Ziffer 3 ungeändert bliebe, so wäre $g_2$ ein
Theiler des Grades der symmetrischen Gruppe von $n-4$ Ziffern, also ein Theiler
von $\Pi(n-4)$. Es müsste also

\[ \frac{n\,\Pi(n-4)}{\Pi(n-3)} = \frac{n}{n-3} \]

eine ganze Zahl sein, also $n \geqq 2n - 6$ oder $n \leqq 6$.

Aus $\varepsilon$) schliessen wir, dass es in $Q$ eine Permutation $\varkappa$
giebt, durch die $0, 1, 2, 3$ in $0, 1, 2, 4$ übergeht, worin 4 eine von 3
verschiedene Ziffer ist. Nach $\beta$) enthält also $Q$ auch die Permutation:

\[ \varkappa^{-1}(0,1)(2,3)\varkappa = (0,1)(2,4), \]

folglich auch:

\[ (0,1)(2,3)(0,1)(2,4) = (2,3,4), \]

also einen dreigliedrigen Cyklus, was der Voraussetzung widerspricht. Hiernach
ist das Theorem II. vollständig bewiesen.

Dass die Fälle $n = 4$ und $n = 6$ wirklich Ausnahmen bilden, geht aus Bd. I,
\S.\,160 und \S.\,182 hervor, wo wir gesehen haben, \gesperrt{dass die symmetrische
Permutationsgruppe von vier Ziffern einen transitiven Theiler vom Index 3 und
die von sechs Ziffern einen transitiven Theiler vom Index 6 besitzt.}

\begin{center}
\rule{3cm}{0.4pt}
\end{center}
